% =====================================================================
% SECTION 2 — BACKGROUND AND RELATED WORK
% Scope: specs/secoes.md, "Seção 2"
% Only equations authorized in specs/arquitetura.md Section 3 may appear here.
% No classification-metric equations.
% =====================================================================

\section{Background and Related Work}
\label{sec:fundamentacao}

This section establishes the material the proposed pipeline operates on and the analytical tools it
applies. Section~\ref{subsec:familias} characterizes the biomedical data families by their physical
encoding, which is the property the detection mechanism actually observes.
Sections~\ref{subsec:radiomica} and~\ref{subsec:sinais} review the descriptor families used for
imaging and for physiological signals, respectively. Section~\ref{subsec:ingestao} surveys ingestion
architectures and data contracts, and positions the present contribution against existing work.

\subsection{Biomedical Data Families and Their Encodings}
\label{subsec:familias}

A recurring source of confusion in multimodal pipelines is the conflation of the \emph{clinical
semantics} of a recording with its \emph{physical encoding}. A twelve-lead electrocardiogram is
semantically a temporal physiological signal regardless of whether it is distributed as a WFDB
record, as a European Data Format recording, or as a comma-separated table of sampled amplitudes.
The encoding, not the semantics, determines which library can read the file and which descriptors
can be computed from it. This work therefore organizes the input space by encoding.

Family F1 comprises temporal physiological signals: electrocardiography, electroencephalography,
electromyography, electro-oculography, photoplethysmography, blood pressure and spirometry. Its
encodings include WFDB record and header pairs, European Data Format recordings and their extended
biosemi variant, MATLAB arrays, vendor-specific spirometry documents in XML, and motion-capture and
polysomnography containers.\needcite{PhysioNet / WFDB}\needcite{EDF standard} Family F3 comprises
two-dimensional imaging studies encoded as single DICOM instances, in which pixel data is
accompanied by an acquisition header carrying modality, spatial calibration and display
parameters.\needcite{DICOM standard} Family F4 comprises volumetric acquisitions --- computed
tomography, magnetic resonance and positron emission tomography --- encoded as NIfTI or MetaImage
volumes, or reconstructed from a directory holding a DICOM series.\needcite{NIfTI format}
Tabular clinical records are encoded as delimited text, and plain two-dimensional images as
conventional raster formats without an acquisition header.

The distinction matters for a further reason: the amount of contextual metadata differs sharply
across encodings. A DICOM instance carries its own modality declaration and spatial calibration; a
NIfTI volume carries voxel spacing but little acquisition context; a delimited text file carries
nothing beyond column names. Any modality-agnostic pipeline must therefore tolerate a wide range of
metadata availability rather than assuming a uniform descriptive header.

\subsection{Radiomics and Handcrafted Feature Extraction}
\label{subsec:radiomica}

Radiomics denotes the extraction of large numbers of quantitative descriptors from medical images,
under the hypothesis that image texture and shape encode clinically relevant information not
apparent to visual inspection.\needcite{radiomics foundational work} A substantial body of work has
standardized the definition and computation of such descriptors so that values obtained by
different implementations remain comparable.\needcite{image biomarker standardization}

Two broad strategies exist for producing a numerical representation of an image. Learned
representations, typically obtained from convolutional networks, require a training stage and large
labeled corpora, and yield features whose individual meaning is not directly interpretable.
Handcrafted descriptors are computed analytically from explicit definitions, require no training,
and remain interpretable in clinical terms. The present work adopts handcrafted descriptors at the
ingestion stage for three reasons: they impose no training dependency on a stage that must run
before any model exists; they behave predictably in the small-sample regimes common in clinical
datasets; and their interpretability is a prerequisite for the auditability the system targets. A
learned-representation branch is retained as an extension point but is not evaluated here.

Second-order texture is characterized through the gray-level co-occurrence matrix
(GLCM).\needcite{Haralick texture features} Let $p(i,j)$ denote the normalized probability of
observing gray levels $i$ and $j$ at a given offset. Entropy, which increases with tissue
heterogeneity, is defined as
\begin{equation}
    H = -\sum_{i,j} p(i,j)\,\log p(i,j),
    \label{eq:entropia}
\end{equation}
and contrast, which increases with local intensity variation, as
\begin{equation}
    C = \sum_{i,j} (i-j)^{2}\, p(i,j).
    \label{eq:contraste}
\end{equation}
Homogeneity, $\sum_{i,j} p(i,j) / (1 + |i-j|)$, and energy, $\sum_{i,j} p(i,j)^{2}$, complete the
texture set adopted in this work.

Morphological regularity is quantified by circularity, $4\pi A / P^{2}$, and by solidity,
$A / A_{\text{hull}}$, where $A$ and $P$ denote the area and perimeter of the segmented region and
$A_{\text{hull}}$ the area of its convex hull. Both decrease as contours become irregular. For
volumetric data the analogous descriptor is sphericity,
\begin{equation}
    \Psi = \frac{\pi^{1/3}\,(6V)^{2/3}}{S},
    \label{eq:esfericidade}
\end{equation}
with $V$ and $S$ the volume and surface area of the candidate lesion; $\Psi$ attains unity for a
perfect sphere and decreases with morphological complexity.

\subsection{Physiological Signal Processing}
\label{subsec:sinais}

Descriptors for family F1 are drawn from two complementary domains. In the time domain, signal
energy is summarized by the root mean square,
\begin{equation}
    \mathrm{RMS} = \sqrt{\frac{1}{N}\sum_{i=1}^{N} x_i^{2}},
    \label{eq:rms}
\end{equation}
complemented by dispersion, skewness, kurtosis, peak-to-peak amplitude and a signal-to-noise ratio
expressed in decibels. In the frequency domain, the power spectral density is estimated by averaging
periodograms over overlapping windowed segments, a procedure that trades spectral resolution for
variance reduction and is therefore well suited to non-stationary physiological
recordings.\needcite{Welch periodogram averaging} From the resulting spectrum the pipeline derives
total power, the spectral centroid, the dominant frequency, the median frequency and the power
contained in fixed frequency bands.

Modality-specific descriptors are layered on top of this common basis. For electroencephalography,
relative band power is computed as
\begin{equation}
    P_b = \left. \int_{f_1}^{f_2} S(f)\,df \middle/ \int_{f_{\min}}^{f_{\max}} S(f)\,df \right.,
    \label{eq:banda}
\end{equation}
over the delta, theta, alpha, beta and gamma bands. For electrocardiography, R-peak detection yields
the sequence of interbeat intervals from which mean heart rate and heart-rate variability are
derived. Variability is characterized by the root mean square of successive differences,
\begin{equation}
    \mathrm{RMSSD} = \sqrt{\frac{1}{N-1}\sum_{i=1}^{N-1}\left(RR_{i+1}-RR_{i}\right)^{2}},
    \label{eq:rmssd}
\end{equation}
by the standard deviation of the intervals, and by the proportion of successive intervals differing
by more than fifty milliseconds.\needcite{HRV measurement standards} For electromyography, the
envelope root mean square and the median frequency are used. For spirometry, the pipeline recovers
forced vital capacity, forced expiratory volume in one second, their ratio, and peak expiratory
flow.

\subsection{Ingestion Architectures and Data Contracts}
\label{subsec:ingestao}

Interoperability standards define how biomedical data is exchanged, and are mature for both imaging
and clinical records.\needcite{DICOM standard}\needcite{HL7 FHIR} They specify encoding and
transport, however, and not the analytical representation a machine-learning pipeline consumes. The
translation from an interoperable encoding to a feature vector remains an application concern, and
is where per-modality specialization typically accumulates.

Within machine-learning engineering, the notion of a \emph{data contract} --- an explicit, validated
interface between the producer and the consumer of a dataset --- has been advanced as a way to
prevent silent schema drift from propagating into models.\needcite{data contracts / schema
validation in ML systems} Related work applies declarative schema validation at pipeline boundaries
so that malformed input fails early and diagnosably rather than deep inside a numerical
routine.\needcite{data validation in production ML}

Three lines of existing work are adjacent to the present contribution. Single-modality clinical
pipelines achieve high specialization but treat the modality as a compile-time constant.
General-purpose radiomics toolkits offer standardized and extensively validated descriptor
implementations, but assume that the modality and the region of interest are known before extraction
begins, and therefore address the stage \emph{after} the one studied
here.\needcite{radiomics toolkit} End-to-end deep pipelines bypass explicit feature engineering
altogether, at the cost of the interpretability and the training-free operation that the ingestion
stage requires.

What is missing across these lines is a mechanism that treats the structure of the input as
something to be \emph{inferred} rather than declared, and that converts every inferred family into a
single representation before extraction begins. The architecture described in
Section~\ref{sec:metodologia} occupies precisely that gap.
